\documentclass[11pt]{article}
% ======================================================================
%  weekpreamble.tex — shared preamble for the weekly course summaries (EN)
%  Concise, readable, intuition-first. Same box style as the main doc.
%  Compiles with pdflatex.
% ======================================================================
\usepackage[utf8]{inputenc}
\usepackage[T1]{fontenc}
\usepackage[english]{babel}
\usepackage[a4paper,margin=2.2cm]{geometry}
\usepackage{amsmath,amssymb,amsthm,mathtools}
\usepackage{bm}
\usepackage{enumitem}
\usepackage{xcolor}
\usepackage[most]{tcolorbox}
\usepackage{microtype}

\emergencystretch=3em
\allowdisplaybreaks[1]
\setlength{\parindent}{0pt}
\setlength{\parskip}{0.45em}

% ---------- math macros (same names as the main document) ----------
\newcommand{\E}{\mathbb{E}}
\DeclareMathOperator{\Var}{Var}
\DeclareMathOperator{\Cov}{Cov}
\DeclareMathOperator{\Corr}{Corr}
\DeclareMathOperator*{\argmin}{arg\,min}
\DeclareMathOperator{\MSE}{MSE}
\DeclareMathOperator{\trace}{tr}
\newcommand{\Reals}{\mathbb{R}}
\newcommand{\Ints}{\mathbb{Z}}
\newcommand{\Comp}{\mathbb{C}}
\newcommand{\Nat}{\mathbb{N}}
\newcommand{\Prob}{\mathbb{P}}
\newcommand{\Tr}{^{\mathsf{T}}}
\newcommand{\Herm}{^{\mathsf{H}}}
\newcommand{\conj}[1]{\overline{#1}}
\newcommand{\abs}[1]{\left\lvert #1 \right\rvert}
\newcommand{\norm}[1]{\left\lVert #1 \right\rVert}
\newcommand{\ind}{\mathbf{1}}
\newcommand{\eps}{\varepsilon}
\newcommand{\diff}{\nabla}
\newcommand{\acvs}{\gamma}
\newcommand{\acf}{\rho}
\newcommand{\pacf}{\alpha}
\newcommand{\sdf}{\mathcal{S}}
\newcommand{\filt}{\mathcal{L}}
\newcommand{\Bsh}{B}
\newcommand{\ms}{\;\overset{\mathrm{ms}}{=}\;}
\newcommand{\toprob}{\xrightarrow{P}}
\newcommand{\todist}{\xrightarrow{d}}
\newcommand{\iid}{\overset{\text{iid}}{\sim}}

\setlist[itemize]{leftmargin=1.5em,topsep=2pt,itemsep=2pt}
\setlist[enumerate]{leftmargin=1.7em,topsep=2pt,itemsep=2pt}

% ---------- compact, titled (unnumbered) boxes ----------
\tcbset{wkbase/.style={enhanced,breakable,fonttitle=\bfseries,coltitle=white,
  boxrule=0.6pt,arc=1.2mm,left=2.2mm,right=2.2mm,top=1.1mm,bottom=1.1mm,
  before skip=6pt,after skip=6pt,
  attach boxed title to top left={yshift=-3.2mm,xshift=2.4mm},
  boxed title style={boxrule=0pt,arc=0.5mm}}}

\newtcolorbox{defn}[1]{wkbase, colback=blue!4!white, colframe=blue!58!black,
  colbacktitle=blue!58!black, title={Definition --- #1}, top=3mm}
\newtcolorbox{result}[1]{wkbase, colback=red!4!white, colframe=red!60!black,
  colbacktitle=red!60!black, title={#1}, top=3mm}
\newtcolorbox{intu}[1][Intuition]{wkbase, colback=yellow!12!white,
  colframe=orange!72!black, colbacktitle=orange!72!black, coltitle=white,
  title={#1}, top=3mm}
\newtcolorbox{retenir}[1][Key takeaways --- the week's reflexes]{wkbase,
  colback=green!5!white, colframe=green!48!black, colbacktitle=green!48!black,
  title={#1}, top=3mm}
\newtcolorbox{exmpl}[1][Worked example]{wkbase, colback=black!3!white,
  colframe=black!55!white, colbacktitle=black!55!white, title={#1}, top=3mm}

% ---------- week header ----------
\newcommand{\weekheader}[4]{%
  {\small\sffamily\bfseries\color{black!60} TIME SERIES~~$\cdot$~~MATH-342, EPFL}\par
  \vspace{2pt}
  {\LARGE\bfseries Week #1 --- #3\par}
  \vspace{1pt}
  {\itshape\color{black!55} #2}\par
  \vspace{6pt}
  \begin{tcolorbox}[enhanced,colback=blue!3!white,colframe=blue!45!black,
    arc=1.4mm,boxrule=0.7pt,left=3mm,right=3mm,top=2mm,bottom=2mm,before skip=2pt,after skip=10pt]
    \textbf{The big picture.}~ #4
  \end{tcolorbox}
}

\begin{document}
\weekheader{10}{30 April 2026}{Forecasting}{The best predictor is the conditional expectation; for causal ARMA we forecast via the $\psi$-weights (future innovations $=0$) and the mean-square error grows with the horizon.}

The whole point of fitting a model is to predict the future. We want $X_{n+h}$ from the
observed data $X_1,\dots,X_n$: $h>0$ is the \textbf{horizon} (or \textbf{lead}), $n$ is the
\textbf{origin}. Throughout we pretend the model \emph{and} its parameters are \textbf{known},
so the only randomness left is that of the future innovations. The recurring message: optimal
prediction needs nothing but the covariance structure $\{\acvs_\tau\}$, and for a causal ARMA
it is computed by writing the future and \emph{erasing the future innovations}.

\section*{Key definitions}

\begin{defn}{Prediction mean square error (PMSE)}
For any predictor $g(X_1,\dots,X_n)$ of $X_{n+h}$, its quality is
\[
P^{\,n}_{n+h}:=\E\!\left[\big(X_{n+h}-g(X_1,\dots,X_n)\big)^2\right].
\]
Smaller is better; minimising it over a class of predictors defines the ``best'' one in that class.
\end{defn}

\begin{defn}{Best linear predictor (BLP)}
The best \emph{linear} predictor is the linear combination
\[
X^{\,n}_{n+h}=P_{X_1,\dots,X_n}(X_{n+h})=\sum_{j=1}^{n}\beta_j X_j
\]
minimising the PMSE. Here $P_{X_1,\dots,X_n}(\cdot)$ is the orthogonal projection onto the
span of the data. For a \textbf{Gaussian} series the conditional expectation is already linear,
so the BLP \emph{is} the overall best predictor; in general $\E[X_{n+h}\mid\cdot]$ is nonlinear and
we restrict to linear predictors, which need only second-order information.
\end{defn}

\begin{defn}{$h$-step forecast error}
$e_n(h):=X_{n+h}-X^{\,n}_{n+h}$. It is unbiased, $\E[e_n(h)]=0$, and its variance is exactly the
PMSE: $\Var(e_n(h))=P^{\,n}_{n+h}$. For a causal ARMA, $e_n(h)=\sum_{j=0}^{h-1}\psi_j\eps_{n+h-j}$.
\end{defn}

\begin{defn}{Three sources of uncertainty}
A forecast carries (1) \textbf{model} uncertainty (wrong model form), (2) \textbf{estimation}
uncertainty (parameters estimated, not known), (3) \textbf{innovation} uncertainty (random future
$\eps$'s). We assume the model is known and parameters are exact, killing (1)\,\&\,(2) and leaving
only innovation uncertainty. (Bootstrap can restore the others; not covered.)
\end{defn}

\begin{defn}{Gaussian prediction interval}
For a Gaussian process $e_n(h)\sim\mathcal N(0,\Var(e_n(h)))$, so a $(1-\alpha)$ interval for
$X_{n+h}$ is
\[
X^{\,n}_{n+h}\;\pm\;z_{1-\alpha/2}\,\sqrt{\Var\big(e_n(h)\big)},
\qquad z_{0.975}\approx 1.96 .
\]
\end{defn}

\section*{Key results \& formulas}

\begin{result}{Lemma --- the conditional expectation is optimal}
$g=\E[X_{n+h}\mid X_1,\dots,X_n]$ minimises the PMSE over \emph{all} predictors.\\
\textit{Idea:} for any $Y$, $\E[(Y-c)^2]=\Var(Y)+(\E[Y]-c)^2$ is minimised at $c=\E[Y]$; apply this
conditionally and take the tower expectation. This is the gold standard; for Gaussian data it is
linear, motivating the BLP in general.
\end{result}

\begin{result}{Theorem --- prediction equations (zero-mean stationary)}
The BLP coefficients $\beta$ solve the orthogonality conditions
$\E[(X_{n+h}-X^{\,n}_{n+h})X_k]=0$ for $k=1,\dots,n$, i.e.\ \emph{the error is uncorrelated with
every observation}. In matrix form $\Gamma_n\beta=\acvs_{[h]}$ with
$\Gamma_n=[\acvs_{i-j}]$ (Toeplitz) and
$\acvs_{[h]}=(\acvs_{n+h-1},\dots,\acvs_{h})\Tr$. If $\Gamma_n$ is invertible,
\[
\begin{aligned}
X^{\,n}_{n+h}&=\acvs_{[h]}\Tr\,\Gamma_n^{-1}\,X,
&P^{\,n}_{n+h}&=\acvs_0-\acvs_{[h]}\Tr\,\Gamma_n^{-1}\,\acvs_{[h]} .
\end{aligned}
\]
Valid for \emph{any} stationary process; the BLP is always unique even if $\Gamma_n$ is singular.
With mean $\mu$, just centre: $X^{\,n}_{n+h}=\mu+\acvs_{[h]}\Tr\Gamma_n^{-1}(X-\mu\mathbf 1_n)$,
same PMSE.
\end{result}

\begin{result}{Theorem --- causal ARMA forecast via $\psi$-weights}
For a causal ARMA, $X_t=\sum_{j\ge0}\psi_j\eps_{t-j}$, the BLP and its PMSE are
\[
\begin{aligned}
X^{\,n}_{n+h}&=\sum_{j=h}^{\infty}\psi_j\,\eps_{n+h-j}=\psi_h\eps_n+\psi_{h+1}\eps_{n-1}+\cdots,
&P^{\,n}_{n+h}&=\sigma^2\sum_{j=0}^{h-1}\psi_j^2 .
\end{aligned}
\]
\textit{Idea:} every past-based predictor is a combination of $\eps_s,\,s\le n$; matching it to the
observable part of $X_{n+h}$ kills all controllable error, leaving the $h$ future innovations.
\textbf{Practical rule:} write $X_{n+h}=\sum_j\psi_j\eps_{n+h-j}$ and set
$\eps_{n+1},\dots,\eps_{n+h}=0$.
\end{result}

\begin{result}{Proposition --- limiting behaviour}
For causal ARMA, $\psi_j\to0$ exponentially, so as $h\to\infty$:
\[
X^{\,n}_{n+h}\longrightarrow\mu,
\qquad
\Var(e_n(h))=\sigma^2\sum_{j=0}^{h-1}\psi_j^2\longrightarrow\sigma^2\sum_{j\ge0}\psi_j^2=\acvs_0 .
\]
The far future is just the unconditional distribution: forecast $=$ mean, variance $=$ marginal
variance.
\end{result}

\begin{result}{Proposition --- integrated (ARIMA) forecasts}
For ARIMA$(p,d,q)$, $Z_t=\diff^{\,d}X_t$ is ARMA$(p,q)$; forecast $Z$, then \emph{re-integrate}:
$\diff^{\,d}X^{\,n}_{n+h}=Z^{\,n}_{n+h}$. For $d=1$,
\[
X^{\,n}_{n+h}=X_n+\sum_{j=0}^{h-1}Z^{\,n}_{n+h-j}.
\]
If $\E[Z_t]=\alpha\ne0$ then $X^{\,n}_{n+h}\approx X_n+\alpha h$ (a drift line). Crucially the
PMSE \textbf{diverges}, $P^{\,n}_{n+h}\to\infty$: no mean reversion.
\end{result}

\medskip
\textbf{Truncated recipe (the hand method, ARMA \& ARIMA).}
(1) Isolate $y_t$ on the left, everything else on the right. (2) Replace $t\to n+h$. (3) On the
right substitute: future obs $\to$ their forecasts, future errors $\to 0$, past errors $\to$
residuals $\hat\eps_s$ (pre-sample $\to0$). Run $h=1,2,\dots$ in order; each forecast feeds the
next.

\section*{Intuition}

\begin{intu}[Why the interval widens, then saturates]
A point forecast is useless without its spread. The error $e_n(h)$ accumulates the unknown
$\eps_{n+1},\dots,\eps_{n+h}$, so the interval grows with $h$. For \emph{stationary} ARMA the
growth saturates: as $h\to\infty$, $\Var(e_n(h))\to\acvs_0$, so the band converges to
``mean $\pm$ a few marginal standard deviations'' --- forecasting far ahead beats nothing. For
\emph{integrated} (ARIMA) processes it grows without bound.
\end{intu}

\begin{intu}[``Write the future, erase the future innovations'']
The two ARMA theorems are one projection seen two ways. Prediction equations project onto
$\mathrm{span}\{X_1,\dots,X_n\}$ using the ACVS; the $\psi$-weight formula projects onto the
\emph{innovation} space $\mathrm{span}\{\eps_s:s\le n\}$. For invertible ARMA these spans coincide,
so answers agree --- but the innovation view is far easier: take the MA$(\infty)$ form and drop every
$\eps$ you have not yet seen.
\end{intu}

\begin{intu}[Mean reversion vs.\ memory length]
Stationary models forget: an AR(1) gives $X^{\,n}_{n+h}=\phi^h X_n\to0$, decaying smoothly to the
mean. An MA$(q)$ has memory only $q$: for $h>q$ the forecast jumps straight to $\mu$. Integrated
models never forget --- a random walk with drift forecasts the drift line $X_n+\alpha h$ forever.
\end{intu}

\begin{intu}[AR$(p)$ shortcut]
For an AR$(p)$ with $n\ge p$, the one-step forecast is simply
$X^{\,n}_{n+1}=\sum_{i=1}^p\phi_i X_{n+1-i}$ --- no matrix to invert. Reason: subtracting the AR
recursion leaves $\eps_{n+1}$, which is uncorrelated with the past, so the prediction equations are
satisfied automatically.
\end{intu}

\begin{exmpl}[MA(2) one- to three-step forecast]
Let $X_t=\eps_t-\theta_1\eps_{t-1}-\theta_2\eps_{t-2}$, $\eps_t\iid\mathcal N(0,\sigma^2)$. Here
$\psi_0=1,\psi_1=-\theta_1,\psi_2=-\theta_2$, $\psi_j=0$ for $j\ge3$. Zeroing future innovations:
\[
X^{\,n}_{n+1}=-\theta_1\hat\eps_n-\theta_2\hat\eps_{n-1},\quad
X^{\,n}_{n+2}=-\theta_2\hat\eps_n,\quad
X^{\,n}_{n+3}=0\ (=\mu).
\]
Variances $\Var(e_n(h))=\sigma^2,\ \sigma^2(1+\theta_1^2),\ \sigma^2(1+\theta_1^2+\theta_2^2)=\acvs_0$.
A 95\% interval for $X_{n+2}$: $-\theta_2\hat\eps_n\pm1.96\,\sigma\sqrt{1+\theta_1^2}$.
\end{exmpl}

\section*{Key takeaways}

\begin{retenir}
\begin{itemize}
  \item \textbf{Best predictor} $=\E[X_{n+h}\mid\text{data}]$; for Gaussian data this is linear, so
        it equals the BLP. Linear prediction needs only the ACVS.
  \item \textbf{Solve} $\Gamma_n\beta=\acvs_{[h]}$ for the general BLP; centre by $\mu$ if the mean
        is nonzero. The error is orthogonal to every observation.
  \item \textbf{Causal ARMA:} write $X_{n+h}=\sum_j\psi_j\eps_{n+h-j}$ and set future
        $\eps\to0$. PMSE $=\sigma^2\sum_{j=0}^{h-1}\psi_j^2$; get $\psi$-weights from
        $\Theta(B)/\Phi(B)$.
  \item \textbf{Limits:} stationary $\Rightarrow X^{\,n}_{n+h}\to\mu$ and
        $\Var(e_n(h))\to\acvs_0$ (band saturates). ARIMA $\Rightarrow$ drift line and PMSE
        $\to\infty$ (band fans out).
  \item \textbf{ARIMA:} difference $\to$ forecast the ARMA $\to$ re-integrate; or apply the
        truncated recipe to the expanded unit-root equation directly.
  \item \textbf{Interval:} $X^{\,n}_{n+h}\pm 1.96\,\sqrt{\Var(e_n(h))}$ for 95\% (Gaussian).
\end{itemize}
\end{retenir}

\end{document}
